\section{Conclusion}
\label{sec:conclusion}

Cloud-native engineering solved restoration and called it recovery. The distinction between the two---long understood by practitioners, long absent from our formal vocabulary---turns out to be exactly the kind of gap a model can close: what separates a restored system from a recovered service is knowledge (recovery-specific dependencies, order, gates, credentials, validation criteria) and substantiation (whether anyone has recently proved any of it works), and both are representable.

The Recovery Graph represents them: a typed graph over services and recovery resources whose edges carry recovery semantics rather than runtime observations, whose vertices move through capability levels that make \emph{restored} a formal waypoint rather than a triumphant endpoint, and whose claims carry expiring, provenance-linked evidence. From this structure follow computable recovery paths and schedules with lower-bound guarantees, bottleneck and counterfactual analyses, circular-bootstrap detection at authoring time, six metrics with stated properties and stated failure modes, and a direct embedding into standard graph query languages. The computed feasibility illustration shows the machinery working end-to-end---and, in passing, shows three quarters of a real architecture's recovery dependencies to be invisible to runtime observability, a single stale credential collapsing the substantiated confidence of an entire tier-1 chain, and a bootstrap cycle caught by a consistency check instead of an outage.

What the paper does not show is effect: whether real organisations can afford the modelling, how far recovery structure diverges from runtime structure in the wild, whether graph-derived plans beat expert runbooks, and whether queryability changes what operators and auditors can do. These are exactly the four questions the registered evaluation design and the open artifact are built to answer, and we invite the community to answer them adversarially. If the Recovery Graph survives that contact, operational continuity gains what infrastructure gained a decade ago: a declarative, versioned, testable representation of intent---this time, for the way back up.
